\documentclass[../Main.tex]{subfiles}


\begin{document}

\section{Theoretischer Hintergrund \& Methodik}

\subsection{Definition des Hack-and-Slash-Genres} \label{sec:Genredefinition}
Ein Videospielgenre zu definieren ist nicht immer einfach, denn Genres sind nicht statisch, sondern entwickeln sich laufend weiter und verzweigen sich. Zudem existiert selten eine scharfe Definition; vielmehr ergibt sich ein Genre aus den Gemeinsamkeiten vieler Spiele, wobei nicht jedes Spiel alle Merkmale erfüllen muss. Im Folgenden beschreiben wir deshalb unser Verständnis des Hack-and-Slash-Genres als Arbeitsdefinition für diese Arbeit. In der Videospielindustrie bezeichnet \textit{Hack-and-Slash} ein Subgenre der Action Games.\autocite{wikipedia_hackandslash} Der Fokus liegt dabei auf Echtzeitkampf mit Nahkampfwaffen, wie Schwertern, Äxten oder anderen Klingen. Fernkampfwaffen sind zwar vertreten, spielen aber nur eine untergeordnete und unterstützende Rolle. Das Genre entstammt den 2D-Beat'em-Up-Games, entwickelte sich weiter in der 3D-Ära und wurde vor allem durch Titel wie \game{Devil May Cry} (2001)\autocite{wikipedia_devilmaycry} oder \game{Bayonetta} (2009)\autocite{wikipedia_bayonetta} geprägt.

Die modernen Spiele zeichnen sich durch hohe Spielgeschwindigkeiten und stilvolle, fliessende Bewegungen aus. Das Verketten präziser Angriffs-Combos, das gezielte Abbrechen von Bewegungen, I-Frames und strategisches Distanzmanagement zum Gegner zeichnen das Gameplay aus. Das, was ein überzeugendes Hack-and-Slash ausmacht, ist nicht allein die Anzahl dieser Mechaniken, sondern das Zusammenspiel dieser. Erst wenn präzises Kampfverhalten und befriedigendes Spielgefühl aufeinander abgestimmt sind, entsteht das genredefinierende Gefühl von Kontrolle und Ausdrucksstärke.\autocite[s.5-s.7]{lindgren2025striking}

Innerhalb des Genres hat sich für Spiele, die diese Ausdrucksstärke ins Zentrum stellen, die Bezeichnung \textit{Stylish-Action} etabliert: Titel wie die bereits genannten \game{Devil May Cry} und \game{Bayonetta} fordern den Spieler auf, nicht nur effektiv, sondern möglichst elegant und abwechslungsreich zu kämpfen, und bewerten dies teils explizit mit Stil-Wertungen im Kampf. \game{Project-Polearm} verortet sich in dieser Tradition.

\game{Project-Polearm} ist ein \textit{lineares} Hack-and-Slash und legt den Fokus besonders auf \textit{fast-paced combat}, also auf einen schnellen Kampffluss. Die Spieler sollen für aggressives Verhalten und stylischen Combat belohnt werden. Dabei soll sich die Spielfigur nicht schwerfällig, sondern agil durch die Level der Spielwelt bewegen.

\subsection{Das MDA-Framework}
Um das Zusammenspiel von Gegnern und Spielfigur systematisch zu erfassen und vergleichbar zu machen, bedarf es eines geeigneten Analyserahmens. Im folgenden Abschnitt wird das MDA-Framework als methodisches Werkzeug eingeführt, das die Beziehung zwischen Spielregeln, emergenten Verhaltensmustern und dem emotionalen Spielerleben strukturiert beschreibt.

Das MDA-Framework (Mechanics, Dynamics, Aesthetics) wurde von Hunicke, LeBlanc und Zubek entwickelt und ist ein formaler Ansatz zur Analyse und zum Verständnis von Gamedesign. Es beschreibt Spiele nicht als statisches Medium wie Buch oder Film, sondern als interaktives Erlebnis, welches aus dem Zusammenspiel dreier klar abgegrenzter Ebenen entsteht. Dieses Verständnis von Spielen als Zusammenspiel mehrerer Ebenen weist Parallelen zu Steven Swinks Konzept des \textit{Game Feel} auf, das er in seinem Werk \game{Game Feel: A Game Designer’s Guide to Virtual Sensation} beschreibt. Swink definiert Game Feel als die subjektive Wahrnehmung der Steuerung und Rückmeldung eines Spiels und untersucht, weshalb sich bestimmte Spielmechaniken für den Spieler präzise, befriedigend oder wirkungsvoll anfühlen.\autocite[s.1-s.6]{swink2008gamefeel} Wenn in dieser Arbeit von \textit{Spielgefühl} die Rede ist, insbesondere in der Erkenntnisfrage, ist damit Game Feel im Sinne Swinks gemeint; die beiden Begriffe werden im Folgenden synonym verwendet.

Die unterste Ebene, die Mechaniken, umfasst die konkreten Regeln und Systeme eines Spiels. Darunter versteht man alles, was durch Daten und Algorithmen getrieben ist, wie Schadenswerte, Treffererkennung, Bewegungsgeschwindigkeit oder die Gegner-KI. Daraus entstehen zur Laufzeit die Dynamiken, also das emergente Verhalten, das sich aus dem Zusammenspiel von Mechaniken und Spielerentscheidungen ergibt. Konkret geht es um den Kampfrhythmus, taktische Ausweichmuster oder eskalierende Gegnerwellen. Die oberste Ebene, die Ästhetik, beschreibt schliesslich die emotionale Reaktion des Spielers auf dieses System. Gefühle wie Kontrolle und Befriedigung bei gelungenen Combos stehen im Vordergrund.

Entscheidend ist die Kausalität zwischen den Ebenen. Mechaniken erzeugen Dynamiken, die wiederum Ästhetik erzeugen. Gleichzeitig erleben die Spieler das Spiel in umgekehrter Reihenfolge. Für die Analyse bestehender Spiele bedeutet das, dass man vom Spielgefühl beginnt und sich systematisch zu den Ursachen, also den Mechaniken, vorarbeiten muss.

Dieses Framework eignet sich besonders gut für die Analyse von Hack-and-Slash-Spielen, da das Genre ein besonders hohes Mass dieser Kausalität aufweist. Beispielsweise sind die genauen Parameter einer einzelnen Mechanik, wie einer Attacke, bedeutungslos, wenn sie nicht im Kontext weiterer Systeme, wie Gegnerlebenspunkten, Gegnerverhalten und Animationsgeschwindigkeit, betrachtet werden.\autocite{hunicke2004mda}

Auf Basis dieses Frameworks wurde ein strukturierter Fragebogen mit 15 Fragen entwickelt, der die drei Ebenen getrennt betrachtet. Er dient als Grundlage für die Analyse der Referenzspiele ab dem Abschnitt~\ref{sec:Referenzspielanalyse} sowie als inhaltliche Basis für den Playtest-Fragebogen.


\subsection{Analyse-Fragebogen} \label{sec:Analysefragebogen}

Die 15 Fragen sind den drei MDA-Ebenen zugeordnet und werden in umgekehrter Reihenfolge gestellt, da die Analyse beim Spielgefühl beginnt und sich systematisch zu den Ursachen vorarbeitet. Thematisch sind die Fragen so ausgerichtet, dass sie gezielt Gegnerverhalten, Gegnerplatzierung, Spielfluss, Schwierigkeit und Feedback beleuchten, da diese Aspekte im Zentrum unserer Erkenntnisfrage stehen.

\subsubsection{Aesthetics (Ästhetik und emotionales Erleben)}

\begin{enumerate}[label=\textbf{A\arabic*}, leftmargin=1.5cm, itemsep=0.5em]
    \item Welche emotionale Reaktion löst der Kampf primär aus: Kontrolle, Befriedigung, Anspannung oder Frustration? Warum?
    \item Wie gut vermittelt das audiovisuelle Feedback die Wucht und Präzision der Treffer? Wie?
    \item Vermittelt die Gegnerplatzierung ein Gefühl von intentionalem Design, wirkt die Anordnung der Gegner durchdacht oder zufällig und wie beeinflusst diese das Empfinden des Spielers?
    \item Verändert sich das emotionale Erleben im Verlauf des Levelabschnitts und ist eine Spannungskurve erkennbar?
    \item Unterscheiden sich die Combos darin, welche Emotionen sie auslösen, bringen sie den Spieler in Bedrängnis oder gibt es ein Gefühl der Erleichterung?
\end{enumerate}

\subsubsection{Dynamics (Spieldynamik und Interaktion)}

\begin{enumerate}[label=\textbf{D\arabic*}, leftmargin=1.5cm, itemsep=0.5em]
    \item Wie sieht der Levelabschnitt von Beginn bis Ende aus, erläutere erkennbare Muster und wie etwaige mit der Spannungskurve zusammenhängen?
    \item Erzeugt die Kombination verschiedener Gegnertypen innerhalb eines Encounters taktischen Druck oder agieren die Gegner unabhängig voneinander?
    \item Entsteht durch Gegnerplatzierung/Spawnpunkte und deren Anzahl ein flüssiger Spielfluss, gibt es einen spürbaren Unterschied zwischen Einzel- und Gruppenkonfrontationen?
    \item Wie erholt sich der Spieler nach einem Fehler, ist die Feedback-Schleife zwischen Fehler und Korrektur nachvollziehbar?
    \item Gibt es verschiedene Combos/Kampfabläufe und folgen daraus verschiedene Arten, die Gegner zu bekämpfen?
\end{enumerate}

\subsubsection{Mechanics (Spielmechaniken)}

\begin{enumerate}[label=\textbf{M\arabic*}, leftmargin=1.5cm, itemsep=0.5em]
    \item Wie unterscheiden sich die Gegnertypen auf Mechanikenebene: Bewegungsgeschwindigkeit, Angriffsreichweite, Schadenswerte?
    \item Wie ist das Angriffsverhalten der Gegner implementiert, reagieren sie aktiv auf den Spieler oder folgen sie festen vordefinierten Mustern?
    \item Wie begrenzt das Leveldesign die Bewegungsoptionen von Spieler und Gegner, gibt es räumliche Einschränkungen, die den Kampf beeinflussen?
    \item Wie sind das Hitstun-System und das Knockback-System gestaltet, lassen sich verschiedene Arten beobachten?
    \item Wie hängen die gedrückten Tasten mit den visuellen Kampfereignissen zusammen, gibt es eine bestimmte Intention hinter den Tasten?
\end{enumerate}


\subsection{Vorgehen bei der Leveldesign-Recherche}
Während die Referenzspielanalyse (Abschnitt~\ref{sec:Analysefragebogen}) den Einfluss des Gegnerverhaltens auf das Spielgefühl am konkreten Spiel untersucht, stützt sich die Recherche zu Leveldesign und Gegnerplatzierung auf eine fokussierte Literaturrecherche, nicht auf eine systematische Analyse. Ziel ist die gezielte Herleitung von Gestaltungsprinzipien für einen zusammenhängenden Kampfabschnitt im gewählten Genre.

Bei den Quellen handelt es sich vor allem um Fachliteratur und Erfahrungswerte von Entwicklern. Sie umfassen zwei Fachbücher, \game{Level Up! The Guide to Great Video Game Design} von Scott Rogers\autocite{rogers2014levelup} und \game{An Architectural Approach to Level Design} von Christopher W. Totten\autocite{totten2014architectural}, sowie vier Fachvorträge der Game Developers Conference (GDC): \enquote{Ten Principles for Good Level Design}\autocite{gdc_tenprinciples2010}, \enquote{Level Design in a Day: A Series of First Steps}\autocite{gdc_leveldesignaday2020}, \enquote{Moving Mountains: Speedy Level Creation as a Desperate Indie Dev}\autocite{gdc_movingmountains2020} sowie \enquote{Level Design Workshop: Balancing Action and RPG in Horizon Zero Dawn Quests}\autocite{gdc_horizonquests2018}. Damit werden bewusst Buchtheorie und praktische Erfahrung verbunden.

Im Fokus stehen besonders jene Themenbereiche, die für den Prototyp relevant sind. Da \game{Project-Polearm} nicht auf Exploration oder eine offene Spielwelt zielt, sondern kampfzentriert ist, konzentriert sich die Recherche auf Themenbereiche wie Encounter-Design, räumliche Typologien, Pacing und Spannungskurve sowie Scaffolding. Die Quellen sind häufig allgemein gehalten und nicht auf ein bestimmtes Genre bezogen, weshalb wir sie in Bezug auf das Hack-and-Slash-Genre interpretiert haben. Auf diese Weise lassen sich in Kapitel~\ref{sec:Leveldesign} konkrete und dadurch praktische Designprinzipien extrahieren.

Es handelt sich überwiegend um Fach- und Praktikerwissen erfahrener Entwickler, nicht um empirische oder wissenschaftliche Studien. Zudem behandelt keine der Quellen das Hack-and-Slash-Genre unmittelbar. Diese genrespezifische Lücke wird durch die eigene Referenzanalyse in Kapitel~\ref{sec:Referenzspielanalyse} geschlossen.

\subsection{Vorgehen bei den Playtests}
Das Ziel des Playtests ist, die Thesen, die wir durch unsere Spieleanalysen und durch die Leveldesign Recherche aufgestellt haben, mit unseren Designzielen, die wir in unserer Leitfrage definiert haben, zu vergleichen. Für die Umsetzung stützen wir uns auf das Buch \enquote{Games User Research} von Anders Drachen, Pejman Mirza Babaei und Lennart E. Nacke.\autocite{drachen2018gamesuserresearch}. Bei Games User Research (GUR) liegt der Fokus auf der Spielerfahrung, was für unsere Ziele genau passt.

Durch die Kombination verschiedener Methoden ist es uns möglich, verschiedene Perspektiven und Ebenen einer Spielerfahrung zu erfassen. So kann die Beobachtung eines Spielers Aufschluss darüber geben, was dieser tatsächlich macht, ein Fragebogen das subjektive Empfinden des Spielers erfassen und ein semistrukturiertes Interview genauer auf Auffälligkeiten im Spielverhalten oder Empfinden eingehen.

Unser Ablauf ist wie folgt: Zuerst werden die Spieltester durch ein Pre Questionnaire charakterisiert. Darauf folgt der Playtest, welcher beobachtet und aufgezeichnet wird, um die Reaktionen der Spieler und ihr Spielverhalten zu dokumentieren. Zum Abschluss gibt es noch ein Post Questionnaire und gegebenenfalls ein semistrukturiertes Interview, um bestimmte Aussagen oder Verhaltensweisen noch etwas genauer zu untersuchen.

Surveys beziehungsweise Questionnaires eignen sich unter anderem zur Erfassung von Spielercharakteristika sowie subjektiven Einstellungen und Erfahrungen. Deshalb haben wir sie für genau diese beiden Zwecke im Pre und Post Questionnaire verwendet.

Im eigentlichen Playtest werden zwei verschiedene Aspekte gleichzeitig beobachtet. Zum einen kann man durch das beobachtete Spielverhalten Rückschlüsse darauf ziehen, ob der Spieler beispielsweise verwirrt ist, wiederholt Fehler macht, auffällige Entscheidungen trifft oder Orientierungsschwierigkeiten hat. Dabei wird zunächst nicht auf die Ursache geschlossen, sondern nur beobachtet und dokumentiert. Während des Spielens werden die Spieler dazu aufgefordert, spontane Gedanken und emotionale Reaktionen laut zu äussern. Dieser Prozess bietet Einblicke in die Wahrnehmung der Spieler. Im Kontext von Hack and Slash Spielen ist eine kontinuierliche Verbalisierung jedoch nicht ganz so einfach, da sehr vieles gleichzeitig passiert und die kognitive Belastung während anspruchsvoller Situationen hoch sein kann. Deshalb werden die Spieler nicht dazu aufgefordert, ihre Gedanken durchgehend zu verbalisieren, sondern vor allem spontane und starke Reaktionen zu äussern. Dadurch können unmittelbare Reaktionen eingefangen und Hinweise auf die Gedanken und Hintergründe von Aktionen der Spieler gewonnen werden.

Das semistrukturierte retrospektive Interview dient dazu, während des Gameplays beobachtete Auffälligkeiten und relevante Situationen nachträglich zu validieren und deren Ursachen aus Sicht der Testperson zu untersuchen.

Teilnehmer Auswertungsmethode kommt noch.

\subsection{Fragebögen für Playtest}

\subsubsection{Pre-Questionnaire}
Das Pre-Questionnaire ist dazu da, um den Playtester zu charakterisieren. Diese Informationen werden dafür benötig, um die Aussagen des Teilnehmers im Kontext zu sehen und für die Gewichtung ihrer aussagen.

\begin{enumerate}
\item Wie viele Stunden spielst du pro Woche Videospiele?
\item Was für Videospielgenres Spielst du besonders gerne?
\item 
\end{enumerate}

\subsubsection{Post-Questionnaire}

\subsubsection{Beobachtungskriterien}

\subsubsection{Interview Fragen}

Teilnehmer Auswertungsmethode kommt noch.


\end{document}